\subsection{Experimental Design}
\label{sec:design}

\paragraph{Data.}
Dataset questions are held out from every stage of construction and only score the final graph. No loop, gate or prompt sees them. The pipeline generates and uses its own curation CQs freely.

The main corpus is 2WikiMultihopQA (2Wiki)~\cite{ho2020twowiki}. A \emph{shard} is a group of questions together with the union of their context paragraphs, with 8 questions in Exp4 and 16 in Exp5, Exp7 and Exp8. From Exp4 on, each 2Wiki experiment excludes the questions used in the earlier experiments. The second corpus is a movie-domain subset of T-REx~\cite{elsahar2018trex} with 482 documents and 14 relation types, of which 386 documents and 507 questions are held out for evaluation. Exp1 uses a bank of building-ontology questions and appears only in Appendix~\ref{app:experiments}. Exp7 holds 480 held-out questions in 30 shards and two run seeds, with no failed cells. On T-REx, the Exp7 arms use Qwen on 33 held-out shards with one-hop gold. Two lists were used there: the T-REx arms of the Exp7 comparisons (Table~\ref{tab:newchecks}, Fig.~\ref{fig:costquality}) used the 2Wiki relation names, and the arms of Table~\ref{tab:trexseed} used the list induced from the T-REx corpus. The 2Wiki checks and Table~\ref{tab:trexseed} use two run seeds; the T-REx Exp7 arms and the AutoSchemaKG arm use one.

\begin{table*}[t]
\centering
\small\setlength{\tabcolsep}{3.5pt}
\caption{The eight studies (Exp6 is an error analysis, the others each vary one component), the Gemma replication and a third model. Each has a plan, a deviation log, code and results in the released archive.}
\label{tab:experiments}
\begin{tabular}{lp{5.7cm}p{5.0cm}l}
\toprule
 & Question & Data & Models \\
\midrule
Exp1 & Does iterative CQ refinement improve a graph? & Building ontology bank, 35 test questions & Qwen \\
Exp2 & Does iteration improve on one batch of the same CQs? & 2Wiki, 24 shards & Qwen \\
Exp3 & Does prompt optimisation improve on a relation list? & 2Wiki, 400 test questions & Qwen \\
Exp4 & Does a CQ $\rightarrow$ SPARQL $\rightarrow$ judge $\rightarrow$ patch loop help? & 2Wiki, 40 shards & Qwen \\
Exp5 & Does a loop with query repair and a rejecting gate help? & 2Wiki, 24 shards $\times$ 2 run seeds & Qwen \\
Exp6 & Which error types occur, and can Shapes Constraint Language (SHACL) see them? & Exp2 graphs & --- \\
Exp7 & What do the combined pipeline's parts add? & 2Wiki, 30 shards $\times$ 2 run seeds & Qwen \\
Exp8 & How do CQ count, budget and hybrid QA interact? & 2Wiki and T-REx, 20 to 33 shards & Qwen, Gemma \\
Replication & Gemma (Exp1, Exp4, Exp5, Exp7, and Exp7's design on T-REx); Exp2 and Exp3 were not repeated & as above & Gemma \\
Third model & Nemotron (Exp5, Exp7 and T-REx; Appendix Table~\ref{tab:replication}) & as above & Nemotron \\
\bottomrule
\end{tabular}
\end{table*}

\paragraph{Exp8 arms.}
The reference arm uses 40 CQs at the 1$\times$ token budget. For each shard the budget comes from a reference run, and an arm stops after using 95\% of it. Other arms are named by CQ count and budget: 160 CQs at 1$\times$, 160 CQs at 4$\times$ and 40 CQs at 4$\times$. The source check is on in the 40-CQ and 160-CQ arms. Each generated CQ is embedded with Qwen3-Embedding-0.6B and compared with the top 5 retrieved chunks, and a CQ that fails is regenerated up to 3 times.

\paragraph{Metrics.}
Evidence recall and all QA measures are scored by code against gold. An LLM judge is used in three places. The builder's own model family (Qwen3.8-Flash-Next) judges the pipeline's internal CQ evaluation (stage 4) and the pass rate of fresh CQs, and a model from the other family judges precision in Exp8. Gemma judges Qwen graphs and Qwen judges Gemma graphs, and the CQ checks use the builder's family.
\begin{itemize}
  \item \emph{Evidence recall} is the share of gold evidence triples present in the graph. A triple counts as present when a chunk-cited triple links the same entity pair under a predicate of the same relation class. \emph{Uncapped recall} counts every such triple. \emph{Capped recall} is recall under a per-question triple budget: the budget is two triples for every source sentence of the question's paragraphs, and each gold evidence counts the probability that a uniform random subset of that question's graph slice no larger than the budget still holds a matching triple, computed exactly as a hypergeometric expectation without sampling, so padding the graph lowers recall. ``Token budget'' always refers to tokens.
  \item \emph{Answer correctness} is scored by code against the gold answer without the grounding requirement: normalised exact match against the gold answer and its aliases or demonyms, the first token for yes/no questions, exact match for dates, and ``unknown'' counts as wrong.
  \item \emph{Grounded QA} asks a model to answer each held-out question from facts retrieved from the graph. Given an input query $q$ and graph $G$, let $R(q, G) \subseteq G$ denote the retrieved supporting fact subgraph. The grounded correctness indicator $\mathbb{I}_{\text{gc}}(q) \in \{0, 1\}$ evaluates whether the answer generated from $R(q, G)$ is correct and grounded in the retrieved facts. Grounded QA over a set $Q$ of held-out questions is $\frac{1}{|Q|}\sum_{q \in Q} \mathbb{I}_{\text{gc}}(q)$. For comparison questions every gold evidence pair (subject and object, in either direction) must appear in a retrieved fact that cites a chunk containing the evidence. For open questions the gold answer must appear as the subject or object of a retrieved fact. The relation name plays no part in grounded QA; relation classes enter only evidence recall (class-aligned recall). QA-time retrieval is the same for every arm and does not use the list: the question is linked to entity labels (with aliases) by exact n-gram matching, the graph is followed for up to two hops, each entity contributes at most 25 facts ranked by BM25, and the facts in $R(q, G)$ are ordered by hop, then BM25, and capped at 60.
  \item The \emph{pass rate of fresh CQs} (Exp7) uses 40 new CQs per shard and run seed. They are written from 40 random chunks with a neutral prompt and no relation list, and no loop uses them.
  \item For each final graph, SPARQL is generated three times for each fresh CQ. A query passes when it returns at most 20 rows, the expected answer span is among them and the judge scores it 2. The pass rate is the mean over the three generations. 
  \item In Exp8, KG-only QA gives the answerer graph facts only. Hybrid QA gives it up to 30 graph facts (main run; a rerun used up to 60) and the top 3 retrieved passages, and \emph{vector QA} gives it the top 5 passages. Measure F2 takes the vector answer unless it abstains, and then the KG-only answer.
  \item \emph{Precision} is judged on 40 sampled facts per graph by the cross-family judge. Gold evidence covers only the facts that questions need, so unaligned triples are not counted as errors. In the fact precision check on the Exp7 graphs, Gemma judges Qwen facts in one pass.
\end{itemize}

\paragraph{Models.}
We implemented this family of CQ-driven methods as one pipeline. Construction, querying and QA use Qwen3.8-Flash-Next (NVFP4 weights) in Exp1--Exp7, and the Gemma replication repeats the experiments named in Table~\ref{tab:experiments} with Gemma-4-26B-A4B-it. NVIDIA Nemotron-3.5-Lightning-30B-A3B, a third model, ran Exp5, Exp7 and T-REx; its graphs stayed at or below .047 grounded QA in every arm, so model comparisons rest on Qwen and Gemma. Exp8 runs Gemma-4-26B-A4B-it and Qwen3.8-Flash-Next on both corpora. All models were served with vLLM, and Qwen ran with thinking off. The QA answerer decodes at temperature 0 (120 output tokens) and extraction at temperature 0 (1,500 output tokens). First-round CQ generation uses temperature 0.7, later CQ steps use 0, and the judges use 0.

\paragraph{External extractors.}
As an external reference, iText2KG 1.1.0 ran with its default settings on each chunk of the Exp7 2Wiki shards with Qwen. It takes no relation list, so none was given, and the Exp7 answerer and scorer scored its graphs. The public extractor EDC~\cite{zhang2024edc} ran on the same chunks, model and engine with its released defaults (Self Canonicalization). It induces and canonicalises its own schema, so it took no list. A third public extractor, AutoSchemaKG~\cite{bai2026autoschemakg}, ran with its released defaults on the same chunks, model family, answerer and scorer, on 30 shards with one run seed; it also took no list.

\paragraph{Registration and statistics.}
Each experiment has a plan, fixed before the main run, that states hypotheses and decision rules; the labels, the checks added later and the run records are in Appendix~\ref{app:registration}. Status labels follow the plan: \emph{registered} hypotheses were fixed beforehand, \emph{secondary} measures are reported beside them, \emph{descriptive} contrasts report intervals only, and \emph{post hoc} analyses and \emph{added arms} came after the results or after the plan was frozen. The unit of analysis is the shard in Exp2, Exp4, Exp5, Exp7 and Exp8 and the question in Exp1 and Exp3; within a shard we average values over run seeds and compare arms pairwise. Intervals are two-sided 95\% bias-corrected and accelerated (BCa) bootstrap intervals that resample by shard, with 10,000 resamples in Exp7 and Exp8. Primary hypotheses are Holm-corrected within the registered primary family of each plan. Tables give the number of shards as $n$ and state units where both appear (30 shards are 60 shard $\times$ run seed units).

\paragraph{Chunk retrieval.}
On the same 2Wiki units of Exp7, two arms answer from the top 5 chunks retrieved by embedding. Top 5 matches the graph arm's prompt length (791 against 794 tokens, median); one arm also receives the 22-name list as a hint. Answer correctness is the measure, since grounding is defined on graph facts.


\paragraph{Controls for volume.}
Tokens per unit count the construction of an arm and its own CQ evaluation: extraction, CQ generation, the pipeline's SPARQL and judge calls and repair, with retries counted and cached calls billed at their nominal cost. They exclude the held-out QA that scores the graph. The construction of an induced list is a separate step and is not included.
Most comparisons hold volume fixed in one of three ways. Arms with equal token budgets stop at the same share of a reference budget (Exp7, Exp8), and the combined pipeline's budget was the token total of LOOP with the list on the same unit. Its arm without the list used the same budget (837,655 and 838,162 tokens). Exp4 and Exp5 compare evaluation-guided patches with random patches of the same size. In Exp7 we also match fact counts. The check at equal fact count draws a uniform random subset of the facts in each stored U-ext graph, with and without the list. The subset matches the fact count $|E|$ of the corresponding L0 graph, and the same answerer and scorer score it. Its plan was fixed before any model call. Two selections reduced the same graphs to the same fact count. One keeps the facts whose relation is most frequent, the other the facts whose end entities are most central. Both read only the facts already in the graph. A chain-preserving and a model selector read no questions, gold or CQs either and keep exactly the guided pass's fact count: the first keeps the facts that add the most connected two-hop entity paths, and the second has the answer model score each fact's informativeness from its source chunk alone. Two more selectors, over the same stored graph and fact count, read no questions and no gold either: a selection by the CQs (BM25 between each fact and the 40 CQs that the guided pass generated for that unit; top facts kept) and a chunk-balanced selection (facts taken round-robin by chunk, ties by hash).
